% ─── Данные и методология ─────────────────────────────────────────────────
\begin{frame}[shrink=10]{Данные и предобработка}

  \begin{columns}[T]
    \column{0.48\textwidth}
    \begin{block}{\faDatabase~Набор данных}
      \small
      \begin{itemize}
        \setlength\itemsep{2pt}
        \item \textbf{Источник:} ЦБ Республики Узбекистан
        \item \textbf{Пара:} USD/UZS (официальный курс)
        \item \textbf{Период:} сентябрь 2017 --- сентябрь 2025
        \item \textbf{Объём:} 2\,087 наблюдений (рабочие дни)
        \item \textbf{Пропуски:} отсутствуют
      \end{itemize}
    \end{block}

    \column{0.48\textwidth}
    \begin{block}{\faCogs~Предобработка}
      \small
      \begin{itemize}
        \setlength\itemsep{2pt}
        \item Отсечение данных до 2017~г. (либерализация рынка)
        \item MinMax-нормализация $[0, 1]$
        \item Формирование лаговых признаков
        \item Разделение: 80\% train / 20\% test
        \item Хронологический порядок (без перемешивания)
      \end{itemize}
    \end{block}
  \end{columns}

  \vspace{0.3cm}
  \begin{tcolorbox}[infobox]
    \small \textbf{Почему с 2017~г.?} До либерализации курс фиксировался административно. Включение этих данных внесло бы структурный сдвиг, искажающий обучение моделей.
  \end{tcolorbox}
\end{frame}
